% =====================================================================
%  R. Nielsen, ON THE CONCEPT OF PRECISION, with Special Reference to Prof. Steen's "Contribution".
%  Postscript to "Philosophy and Mathematics".  Copenhagen: Gyldendal (F. Hegel), 1860.  55 pp.
%  (Om Begrebet Nøiagtighed med særligt Hensyn paa Hr. Prof. Steens »Bidrag«.  Efterskrift til »Philosophie og Mathematik«.)
%  ENGLISH TRANSLATION, from transcription.tex (Danish).
%  GENERATED by ebook-method/begrebet-noiagtighed/tr_build.py --write from this template and the two translators' parts in
%  texts/nielsen/begrebet-noiagtighed/.parts/tr/ (TRANSLATOR-BRIEF.md; terms: GLOSSARY.md, Nøiagtighed = precision
%  approved by Hans 2026-10-10).  Conventions: TRANSLATION-PLAYBOOK.md.  Page numbers of the 1860 print in the margin.
%  Emphasis as in the Danish: \textbf = bold Fraktur, \emph = letter-spacing, \textit = antiqua (Latin, French, English).
% =====================================================================
\documentclass[12pt,a4paper,openany]{book}
\usepackage[utf8]{inputenc}
\usepackage[LGR,T1]{fontenc}
\usepackage{textalpha}
\usepackage{amsmath}
\usepackage{libertinus}
\usepackage{graphicx}
\DeclareUnicodeCharacter{0254}{\reflectbox{c}}  % ɔ (U+0254) = reversed c
\DeclareUnicodeCharacter{A75B}{r}   % r rotunda (ꝛc. = etc.)
\usepackage[english]{babel}
\usepackage[protrusion=true,expansion=true]{microtype}
\usepackage{setspace}
\usepackage[margin=1.3in]{geometry}
\usepackage{fancyhdr}
\usepackage{hyperref}
\hypersetup{pdftitle={On the Concept of Precision}, pdfauthor={Rasmus Nielsen}, hidelinks}
\setlength{\headheight}{15pt}
\pagestyle{fancy}
\fancyhf{}
\fancyhead[LE,RO]{\thepage}
\fancyhead[C]{\textit{R. Nielsen: On the Concept of Precision (1860)}}
\renewcommand{\thefootnote}{\ifcase\value{footnote}\or *)\or **)\or ***)\else
  \arabic{footnote})\fi}
\setstretch{1.3}
\usepackage{marginnote}
\newcommand{\opage}[1]{\marginnote{\footnotesize\textit{[#1]}}}
\providecommand{\tg}{\operatorname{tg}}

\begin{document}

% PDF 5, the title page, p. [1]
\begin{titlepage}
\begin{center}
\vspace*{3em}
{\large On}\\[1em]
{\Huge\textbf{the Concept of Precision}}\\[1em]
{\small with Special Reference to}\\[1em]
{\Large\textbf{Prof. Steen's ``Contribution''.}}\\[1.5em]
\rule{4em}{0.4pt}\\[1.5em]
{\small Postscript to ``Philosophy and Mathematics''}\\[0.8em]
{\small by}\\[0.8em]
{\large\textbf{R. Nielsen.}}\\[5em]
{\small\textbf{Copenhagen.}}\\[0.5em]
{\small Published by Gyldendal's Bookshop (F. Hegel).}\\[0.5em]
{\footnotesize Thiele's Printing House.}\\[0.5em]
\textit{1860.}\\[3em]
{\small Translated from the Danish}
\end{center}
\end{titlepage}

% ===== TEXT (pp. 3--55) =====
% --- p. 3 ---
\opage{3}\setcounter{footnote}{0}%

\begin{center}\textbf{On the Concept of Precision, with Special Reference to}\\
\textbf{Prof. Steen's ``Contribution''. Postscript to ``Philosophy}\\
\textbf{and Mathematics''.}\end{center}

\begin{center}\rule{3em}{0.4pt}\end{center}

In all sciences, and most especially in the most exact
of all sciences, what matters is to be precise; but
the demand for precision can, as is often enough seen in critics
and notices, nevertheless also contain questionable aberrations. For one can be precise in the inessential and imprecise in the essential, precise in pointing out
particulars but imprecise in grasping the whole, precise
in putting forward one's own humble opinions, but highly
imprecise in interpreting and reproducing the point of view one wishes to
contest.

In ``Philosophy and Mathematics'' there are to be found a couple of rather considerable errors of inadvertence and some lesser imprecisions,
which are here to be corrected. In order, however, to make corrections in the correct manner, one must know how to distinguish between
real and alleged errors; for just as it runs counter to the demand of precision to let real and manifest errors stand without correction, so it is, after all, also a
sign of insecurity, an insecurity which the habit of imprecision only too easily brings with it, to let oneself be taken by surprise by
bold protests and then, in such surprise, be misled into
correcting what is in itself correct, and acknowledging errors where there is
no error.

% --- p. 4 ---
\opage{4}\setcounter{footnote}{0}%
In a piece entitled: ``Contribution to a Critique of Prof.
R. Nielsen's ``Philosophy and Mathematics'', Copenhagen,
Reitzel's Publishers, 1859'', Prof. Steen has, certainly not
without self-conquest and a great sense of duty, at last felt
called upon to serve precision and to help us with the
corrections. Before, however, we avail ourselves of the guidance offered, we shall first, so as not to go astray,
see to indicating in general the direction of the road, by taking note of
some signposts, in that we, in other words, take note of
some of the most conspicuous caricature forms of literary precision.

As the first principal form we name the \emph{pettifogging}. To
point out errors is a sign of superiority; to point out
many errors is a sign of much superiority; to point out tremendously many, tremendously great and tremendously gross errors is
therefore a sign of a tremendous superiority: here we have the
pettifogger, the literary market-crier. That the task which
``Philosophy and Mathematics'' has set itself is certainly not among
the easiest, even those who for the rest
condemn the whole attempt will doubtless concede without contradiction. A severe but thorough critique, which would inform
us of what in such an attempt might perhaps be found
usable, and what in it must be regarded as failed, would
then, in proportion to the difficulty of the task set, have
a task not at all easy either. But what is difficult for a
thorough critique is, naturally, not difficult for a pettifogging critique; the two kinds of critique are opposed to one another, so diametrically opposed, that the more difficult the matter
is for the thorough one, the easier it is for the pettifogging one, the more time the former requires for strenuous deliberations, the more comfortably the latter completes its hasty work
. We can illustrate the opposition by an easy example. ``For what readers is the propaedeutic treatise,
Philosophy and Mathematics, actually intended?'' If the treatise is to be regarded as a draft, intended chiefly for
% --- p. 5 ---
\opage{5}\setcounter{footnote}{0}%
use in oral lectures for younger students:
then a scientific reviewer surely sees at first glance that the regard to the explanations an oral lecture can impart is a regard decisive for the judgment, which must be upheld throughout, so that the objections
which might otherwise perhaps be made against the use of
some in themselves ``difficult examples'', fall away of themselves.
If, on the other hand, one in the judgment wishes to take no
regard whatever to any oral lecture, but to let the treatise ``appear as a printed production, which both lays
claim to and must be prepared for literary critique'': then
it follows from this, as a thorough critic immediately concedes, that every appeal to the knowledge of ``the young students'' must be rejected as unwarranted. Otherwise,
when the pettifogger begins to put forward his higgledy-piggledy
of jumbled points of view and heterogeneous yardsticks.
If the exposition is distinct, then it is broad and trivial, and
what is said so generally known that it is nothing
to offer to the experts. If a dialectical splinter sticks out, then
``one allows oneself, until the contrary is proved, to assume
the touched-on proposition to be meaningless.'' If one comes upon
a more obscure passage; it is to be treated in one of two
ways: either one hastily imputes to it an absurd and
``stark mad'' meaning, declares the imputed meaning to be
``condemned to death'', and, since ``a printed production must
be prepared for literary critique,'' naturally will not hear a
word of oral lectures in which an explanation was to be
given to the young students; or, if ``the printed production'', however much one screws and twists the individual
propositions, will not become quite as meaningless as one
might wish; then one chooses the opposite way out, then one concedes that the expert can indeed see what the
author's thought is, but ``the young students?'' yes, ``the
young students!'' --- with a compassionate glance one turns
toward ``the young students'', and has then at once a worried
% --- p. 6 ---
\opage{6}\setcounter{footnote}{0}%
thought in the great principal idea that ``He who comes forward as
a teacher for the young students, who have recently left
school, takes his examples from a sphere unknown to them.''
By thus shaking together objections, as one shakes
bits of glass in a kaleidoscope, the pettifogger serves scientific precision, and the demand of precision makes its
effect: an inquisitive reader peeps into the kaleidoscope, hears
about the imprecise, sees the precise, and is --- convinced.

As the second principal form of the caricature we can adduce the \emph{pedantic}. Precision and punctuality
coincide, after all, in this, that everything in the particular is to be as to
a dot. Precision is therefore the same as punctiliousness; the pettier the punctiliousness, the greater the precision:
here is the pedant, the literary fusspot. The demand of precision is now this, that no proposition ought to be regarded as
truly clear or properly intelligible which is not provided with
so many hooks of caution, conditions, restrictions
and subtle reservations, that even standing alone, torn out
of all context, it can make every misunderstanding impossible; here there becomes something to prick at, something to discover and
something to lament for him whom it especially must please to
lament, to have had to weary the reader over forty pages, in order to set forth the incorrect, the meaningless and
the unclear in the first fourth part of half a page.''

In most forms of the humane productions of literature and science one otherwise knows quite well how to make
a distinction between precision and pedantry; only the out-and-out mathematical punctiliousness seems still to want to have the privilege of being regarded as strict precision. It is said in
a star catalogue: ``these stars are visible throughout January and
February at the stated hours in the northern sky''; a scrupulous reviewer steps forward and
corrects: ``this statement is imprecise; here it must surely be
presupposed that the northern sky throughout January and
% --- p. 7 ---
\opage{7}\setcounter{footnote}{0}%
February is not obscured either by fogs
or by proper clouds; if these assumptions are incorrect, then
the whole rests on a misunderstanding, for which the author himself bears the blame, for one must guess at
his meaning when he does not make it precise;'' Everyone understands that such a correction is in the highest degree pedantic.
But when it is now emphasized in ``Philos. and Mathem.'' that
mathematics too can in its language express both the
possible and the in itself impossible, in that there is added by way of example, and
moreover cursorily, because a heavy-footed description of a
matter so easily grasped would be unnecessary:
``$a = \frac{b}{c}$ is for thought a reality, $b = \pm\sqrt{a^{2} - x^{2}}$
a possibility''\dots{}, and a reviewer then, on behalf of precision, correcting remarks: ``when it is said that
$a = \frac{b}{c}$ is a reality, it is surely presupposed that $b$ and $c$ are rational,
when $b = \pm\sqrt{a^{2} - x^{2}}$, $x$ must surely lie between $+$
and --- $a$; if these assumptions are incorrect, then our protests rest on a misunderstanding, for which the author himself
bears the blame, for one must guess at his meaning,
when he does not make it precise'': might not such a correction
also fall under the category of the pedantic?\footnote{What a contradiction pedantry contains in itself is seen from this,
that the pedantic correction almost always in turn requires a new correction. Should b and c be presupposed rational, in order that $a = \frac{b}{c}$ may
be a reality, the presupposition in such generality is still unclear and imprecise; no, one must, in order to ``make precise''
the expression, expressly presuppose that the rational b and the rational c are of unequal size. If b and c are not of unequal size, then
the proposed restriction is even wholly incorrect, for
$a = \frac{\sqrt{2}}{\sqrt{3-1}}$ is surely just as much a reality as $a = \frac{2}{3-1}$.
This pedantic correction could then in turn occasion a correction etc.}

% --- p. 8 ---
\opage{8}\setcounter{footnote}{0}%
That the pedantic precision must by its own
consequences turn over into the imprecise, the unreliable and
wholly untrustworthy is a matter of course; when a fuss is to be made
of the inessential, the in itself essential must be
overlooked: if one makes everything equally essential, one makes at the same
time everything equally inessential.

The pedant and the pettifogger are kindred spirits; in their
reviewing procedure there is only this difference, that
the pettifogger crushes with racket when the pedant pricks
to death with a needle. In pettifoggery as well as in pedantry
one presupposes, nevertheless, gladly a certain lack of sound ability; % print: saud (misprint)
but if the ability, though only as narrow professional ability, is also present, and with it a distinct notion of how the pettifogging and pedantic
means are best combined and applied, when one wants to get at an
author and have him ``convicted''; then the measure is
full, then the caricature of precision assumes its highest and most
interesting form, namely the caviling!

The caviler is interesting in that he enters into
reflection; for the interesting has always in one way or
another its attraction in a reflection. The pettifogger
is immediate, and his appearance therefore basically without
essential interest; his loud cries, his violent gestures, his flowering, full-blooded coarsenesses and unabashed
personalities may certainly, insofar as they are ``testimonies to the impartiality, the ruthlessness, the zeal and
warmth with which he for the sake of the cause enters into the
discussion,'' now and then be quite amusing; but interesting they are not, actually. To captivate reflection,
the pettifogger would himself have to have reflection; but to be a
pettifogger not much reflection is needed, not much more
than to be a drummer.

That the interesting is conditioned by reflection does
not at all exclude the possibility that there can also be
a reflection which is by no means interesting. Pedantry
% --- p. 9 ---
\opage{9}\setcounter{footnote}{0}%
is not without reflection, and yet pedantry is not the
interesting. When the pedant, as regards ``the countless imprecisions'', corrects a book with the same precision
as that with which ``the Licentiate'' in ``A Student's Adventure''
describes ``the Lægdsgaard Ølsebymagle'', he can, especially when
he sets about ``lamenting that he has had to weary the reader
over forty pages, in order to set forth the incorrect, the meaningless and the unclear in a little piece of half a page,'' surely
in a way become interesting, namely for the reader's
reflection: ``forty pages, so many errors, so many
imprecisions, so many corrections! He who with such
zeal wishes to serve precision must surely, before he
gets the whole book corrected, and with it the author literarily
executed, correct himself to death!'' But one thing is it that pedantry lets itself be wrung of an interesting side; quite
another that it should by its character itself belong
to the interesting. The pedant can become interesting in that
he can become comical; but the grumbling, sallow, reflection-sick pedant is in reality a pitiable and by no means interesting figure.

The caviler, by contrast, who with force and ability
understands how to seize the elements offered him from
lower spheres, and to use them as means for his purpose, the caviler, who knows both what he himself wants and
also what one ought to want in a literature, the caviler, who
puts down a Scientific and has a Passionate in
mind, just as in simple arithmetic one puts down zero and
has nine in mind, yes, the caviler has, which especially betrays itself when he is a little clumsy, always a doubling, and for that reason belongs from all sides to the
interesting: the more one sees through his nature, the
more interesting he becomes, for his nature is, on account
of the doubling, a reflection.

The doubling in which the caviler has his interestingness, % print: Interssante (misprint)
rests, generally seen, on an alternating semblance of
% --- p. 10 ---
\opage{10}\setcounter{footnote}{0}%
Morality and Science. The Scientific is that ``the many
errors are corrected, and corrected with precision;'' the Moral is
that they are corrected on the right side. Since it is now, as we
with the caviler are surely on scientific ground, ``provisionally not so much a matter of the multitude of the errors, as
of their nature, not so much of their quantity,
as of their quality'': the many corrections must surely above
all aim at giving the errors the right quality, and
getting something out of the quality. We choose again a small
example. In ``Phil. and Mathem.'' the relation between
the first and second differential coefficient is elucidated by a couple of
formulas from mechanics: $v = \frac{dx}{dt}$ and $\Phi = \frac{dv}{dt} = \frac{d^{2}x}{dt^{2}}$. An expert
, but otherwise credulous, reader might perhaps
even run through the little passage p. 61,
at the bottom, several times without noticing anything, since he unwittingly thought
the velocity now constant and now variable, according as
the context and the chosen expressions naturally brought it
with them\footnote{Even if the expert, upon unwittingly stumbling on the error,
should become suspicious that the Author was perhaps at sea with
the concepts: he would nevertheless as a man of honor first look around a little
among the Author's other productions, and e.g. consult Phil. Propaed. pp. 164---167, before publicly and
in print attacking him by stamping the error with the mark of gross
ignorance. But of course: such regards apply
only among men of honor; pettifoggers and cavilers have
their own chapter of duties.}. But, cries the pedant, the reader grasps
at once what has happened here, for here a discovery has happened. And
yet it would help but little that an error was discovered, if
the error did not have the right nature, if there was no
proper quality in the error, if the error were not one to
lament. Instead of taking note of the importance of the discovery
and lamenting the error, ``the young students'' would perhaps even
be capable of taking note of the pedant's importance, and amusing themselves
% --- p. 11 ---
\opage{11}\setcounter{footnote}{0}%
with: quel bruit pour une omelette! No, if the discovery is
to count for something, if there is to be substance in the discovery, then
a man is needed here who can use the discovery, put color
on the correction, and lament the error. The discovered error
is then, lamentably enough, no error of inadvertence, no, it is an
error of ignorance, an error of the grossest kind! ``The Author'', who
wishes to ``clarify the concepts'', has namely not known what yet
``every student knows'', that there is a difference between a
constant and a variable velocity; no, that he has not
known, and that is just the lamentable thing; let us all lament:
lamentable velocity, lamentable formulas, lamentable $v$! lamentable
$dv$ divided by dt, lamentable $\Phi$, o lamentable errors! lamentable that such errors can be committed by one who
comes forward as author and teacher; still
more lamentable that they find defenders!''

Therein consists precisely the caviler's interesting doubling,
that he not only knows how to discover the error, but even to
make its quality so lamentable that here there is in earnest cause
to lament over it.

With these signposts before our eyes we shall now, obeying
the laws of precision, endeavor to wander the right
ways of corrections, and thereby try the guidance Prof.
Steen's Contribution in this respect may offer us. Right
in the introductory declaration the ``Contribution'' (p. 4) promises us a good deal,
which, according to the Prof.'s assurance, will be ``of great importance with regard to the judgment of the Author's
standpoint in general.'' Dare we also rely on
the fulfillment of this promise? ``The Author's'' standpoint is a
philosophical one, more precisely: a dialectical one: what ``contributions''
does the Prof.'s ``critique'' now deliver toward an elucidation of this
standpoint? On the occasion of the question of the significance that must be attributed to the concept: approximation, the
``Contribution'' (p. 27) says that ``it is not the intention here to give
any critique of the dialectic that runs through the whole treatise;'' with this the dialectic is then disposed of. What
% --- p. 12 ---
\opage{12}\setcounter{footnote}{0}%
position Prof. Steen's critique takes toward philosophy
and philosophical development of concepts as a whole, the ``Contribution''
(p. 19) has indicated to us in a quite distinct and easily grasped
manner. Namely, on the occasion of ``the problem of infinity, the great
common problem of mathematics and philosophy,'' the following interesting piece of information is imparted here: ``I,''
(Professor Steen) ``believe now certainly that the Author here is on
a wholly incorrect path. There is in the first place nothing
that can be called `the problem of infinity' in mathematics''\dots{} ``But the Author is Professor of Philosophy, my
opinion on this question means something only to me
myself, and it is also only for my own sake that I have
wished to express this different conception\dots{}'' Thus: instead
of engaging in an examination of ``the Author's'' standpoint, the
mathematician, for safety's sake, takes his stand on the standpoint of ``Faith'',
and then by virtue of an ``I believe'' \textit{(credo, ut}
\textit{intelligam)} asserts his ``opinion'' securely and confidently. That
one may call caution, the caution of a truly believing mathematician!
More cautious has perhaps never yet any
``severe critique'', by keeping at a proper distance, avoided contact with the criticized author's standpoint.
``I shall,'' says Prof. Steen (Contr. p. 19) ``in the
present investigation, as well as I am able, enter
into the Author's train of thought, whereby it is then a stroke of luck for me
that he wishes to follow the path marked out by mathematics.''
Certainly it is a stroke of luck to be a mathematician by profession, when
it is a matter of following a path marked out by mathematics;
but in an investigation that concerns philosophy just as much
as mathematics, ``the path marked out by mathematics''
must surely, reasonably, be so chosen that in the
development of concepts it coincides with the path marked out by philosophy; the mathematician would thus overrate his
``luck'', if in such an investigation he overlooked
the philosophical, and ``believed'' himself entitled to strike down
into the merely mathematical.

% --- p. 13 ---
\opage{13}\setcounter{footnote}{0}%
To elucidate how unreasonable and thoughtless a separation of
the mathematical and the philosophical in an investigation like
this must be, we choose from ``Phil. and Mathem.'' the two
examples which in our conviction are the only ones that
with any justification can be reckoned under the category of essential
errors. The first of these examples is found p. 10, and concerns the opposition between the trigonometric line, cos x,
and the corresponding arc $\mathrm{arc}\, (\cos = x)$; the second
is found p. 43, and concerns the manner in which the presupposition and condition for a development of the Taylor series
is put forward.

Wherein, then, essentially seen, does the error in our first
example consist? Instead of hitting the point in which the error lies,
and determining the relation between its mathematical and philosophical nature, Prof. Steen begins with
a fallacy, in that he from this, that all functions are
intellectual acts, wishes to derive that all functions
must be mutually canceling acts; but that is a
rashness. What he here calls ``an incorrect ground'', is
thus his own invention; likewise the incorrect in
what he calls ``the incorrect assertion'' is solely and alone to be
ascribed to his own assertion. Whatever and however many meanings one may, from a purely mathematical point of view,
be led to connect with such expressions
as $\cos x$, $\mathrm{arc}\, (\cos = x)$ etc., no one will yet be able to
deny that $x$ in $\cos x$ according to its concept is an arc, and $x$ in
$\mathrm{arc}\, (\cos = x)$ according to its concept a straight line. If one
now transforms $x$ in both expressions into pure numbers, and lets the original meaning not merely recede into the background, but
fall away, and, mark well, fall away in such a way that there
is not even a glimmer of the original left: will not then the arc
of a pure number be just as logical as the chord of a
pure conscience?

When Prof. St. thus asserts that in
% --- p. 14 ---
\opage{14}\setcounter{footnote}{0}%
the theory of functions $x$ in $\cos x$ is not an arc, $x$ in $\mathrm{arc}\, (\cos = x)$ not a
straight line, but that both are pure numbers, and that this conception, necessary ``for a
philosophical consideration'', is ``the only correct conception
for the theory of functions'': then one feels strange.
Should ``the theory of functions'' now really
demand that the original meaning of $x$ in $\Phi(x)=\cos x$
and in $\Psi(x)=\mathrm{arc}\,(\cos=x)$ be wholly obliterated? That
would have to be when we considered the mathematical functions from the standpoint of --- ``Faith''; for that the assertion put forward by Prof.
Steen, considered from the standpoint of science, is by no means approved by all ``the mathematicians of Europe'', can e.g. be seen from Professor Jürgensen's ``Elements'' p. 18\footnote{Elements of Mathematical Analysis, Copenhagen
1855.}, where it is expressly stated that
``The logarithm is a function of the number, the sine (as line)
a function of the radius and the arc.''

But even if the error complained of were sufficiently
demonstrated, the question of its influence on the
philosophical would still remain. The error, as it
now stands, could, as far as the merely mathematical is concerned, for all that
well have been twenty times greater, and its
influence on the development of concepts still be reckoned as
nothing; for the brief remark in which it occurs
can without harm to the train of thought be wholly dropped.

That in setting up the presuppositions on which the Taylor series
rests, an incorrect remark has crept in, p. 43, concerning the coefficients $A_{0}$, $A_{1}$, $A_{2}\dots$, shall willingly be conceded, just as it shall also be conceded that something
on this occasion has been set up as a condition which,
more closely considered, is not a condition, but a consequence.
% --- p. 15 ---
\opage{15}\setcounter{footnote}{0}%
Meanwhile, however much fuss one may wish to make of these introductory errors of remark, the Taylor series itself is nevertheless
presented in ``Phil. and Mathem.'' from pp. 44---45, to be sure not by directly copying out Ramus
(Algebra and Theory of Functions pp. 99---100), but yet with
such close adherence to and essential agreement
with the manner of proof given by Ramus, that an attack
on ``the merely mathematical'' must here transform itself into an
attack on Ramus. It therefore gives food for thought that
Prof. Steen, while preparing to complain of
``the manner, dreadful for a mathematician, in which Taylor's formula is treated'', must prepare himself
by a little demonstration against Ramus. Certainly it has
something imposing to hear Prof. Steen on this
occasion appeal to ``the mathematicians of Europe'', just as
it also cannot but instill in the reader respect for
the precision, the fine precision, with which Prof. St. sees everything and notes everything and makes use of everything,
when he, after having appealed to ``the mathematicians of Europe'', then once more with a glance at ``Philosophy and
Mathematics'' makes the fine remark that ``The expression:
Taylor's series is transformed into Maclaurin's series'', is insofar not quite happy, as it is thence not entirely
clear that the latter is merely a plain and simple special case
of Taylor's''; but the all too imposing can then again
easily overwhelm, and the many acute hints, the
manifold many and fine remarks easily divert attention from the main point. The question whether
``the Author'' in Philosophy and Mathematics has shown that he
knows anything about the mathematical in the Taylor series
is surely in this connection the main point; if therefore
``the Author'' can with good right in the main point appeal to
Ramus, then Ramus is --- with all possible respect for
``the mathematicians of Europe'' --- with us in the main point still always Ramus.

% --- p. 16 ---
\opage{16}\setcounter{footnote}{0}%
The mathematical thus treated cannot then, however much Prof. St. troubles himself, nevertheless be torn loose from the philosophical. ``The dreadful manner in which
Taylor's formula'', to say it now a second time, ``is treated, and the dreadful effort Prof. has made,
in order by this treatment to become quit of the philosophical altogether,
has by no means led to the desired result. The proposition about the scientific method as that whose infinite mediation is due at once to ``the necessary movement of the subject matter and the free movement of thought'', holds fast; even after
Prof., by giving up ``the necessary movement of the subject matter'',
has split and divided and twisted and turned and unraveled every
jot for so long, until the whole by his ``free movement of thought'' has become pea-straw, critical-mathematical pea-straw: will the proposition on the scientific method, however
eagerly and vehemently Prof. Steen shakes the straw in his
``Contribution'', remarkably enough, not let itself be shaken off.

``It shall only be mentioned'', says Prof. (``Contr.'' p.
39), ``that the development which the Author gives after Ramus is in
no way given with \emph{necessity} by the matter
itself and ``the free movement of thought'', which is best seen from this,
that the proof set up by Ramus is quite peculiar to him\footnote{It shall only be mentioned that the proof ``quite peculiar to Ramus
is already found in Lagrange \textit{(Calcul des fonctions. A Paris}
$1806. p. 10-13$), and thus --- praised be Prof. Steen's
tremendous precision --- is already ``quite peculiar'' to Lagrange.''}, whereas Taylor's formula, which dates from 1715,
has previously appeared and been proved in other ways''. Golden
words! here is something to learn. When, therefore, for one and the
same theorem two, three, perhaps four equally sufficient proofs are given: what does this prove? Why, it proves that
none of these proofs is ``with necessity''. And when
a mathematician sets up a ``proof quite peculiar to him''
for a theorem which over a hundred years ago
% --- p. 17 ---
\opage{17}\setcounter{footnote}{0}%
has been proved by another proof: what does
this prove? Why, it proves, of course, that ``the free movement of thought'' cannot
be united with ``the necessary movement of the subject matter''! Prof.
Steen makes at several places in the ``Contribution'' a noticeable effort that the reader shall be informed of \emph{how}
he still investigates in ``Cognition and Investigation'', and repeats his ``Logic'' in order ``to note in every proof
what is the actual \textit{punctum quæstionis} and what
is the actual \textit{nervus probandi}: here is a logical yield,
here are fruits of his investigation, sound, beautiful, savory fruits!

He who will now but taste one of these fruits
will, if he otherwise has any taste for scientific
proofs, convince himself without difficulty that Prof.
Steen, when he chooses to criticize the merely mathematical
without engaging with the philosophical, may well have his
good and sufficient reasons.

``I shall'', says Prof., ``in the present investigation, as well as I am able, enter into the Author's
train of thought''; we repeat this praiseworthy utterance, because
we everywhere find it in order that he who would criticize
an author enters into his ``train of thought'', and can
only regret that Prof., despite his promise, not merely
altogether goes around ``the Author's'' train of thought in the whole, but
even, as if he feared making himself dependent on any train of thought whatever, has not even wished to \emph{articulate}, but
merely to \emph{divide up} his Contribution. If an articulation were at issue, there would have to be here, as Prof. St. must also be
reminded each time he repeats his ``Logic'', a
ground of articulation, a \textit{fundamentum divisionis}, which from first
to last could determine the articulation; but such a ground of articulation then also at once brings one into dependence on
a train of thought. This dependence one avoids, on the other hand,
when one, instead of articulating, will merely divide, i.e. carve up,
cut into pieces. ``To facilitate the survey, the investigation''
% --- p. 18 ---
\opage{18}\setcounter{footnote}{0}%
--- it says (p. 4) --- ``is divided into several sections''; i.e.
divided as one divides a piece of hasty work, divided, not as
for the sake of a content, but solely for the sake of the material, ``since the
material is so rich'', divided --- into seven contributions, cut into seven
pieces, arbitrarily, as one cuts a sausage.

The only presupposition under which so arbitrary a division and cutting up, interrupting all connection, destroying all train of thought, could
in any measure be justified, would then have to be the circumstance that each of the
set-up ``divisions'' found its center, its natural
point of support, in the treatment of a determinate main error, which
by its great importance had to captivate attention.
We shall now examine to what extent such a circumstance
is present, and then, so as not to skip in the text, begin
with the beginning.

I. ``The Author is (Contr. p. 47) not quite clear about
the relation between the individual elementary functions''; that is the assertion which in the first division (pp. 4---6)
constitutes the theme of the critique, whereby it is then especially emphasized that
``the power'' $x^{a}$ and the exponential function $a^{x}$ are now separated as two independent functions, now again thrown
together into one. The reader who takes the matter seriously,
and does not wish to regard Prof. Steen's ``Contribution to a
Critique'' of a propaedeutic treatise as a contribution to a
characterization of literary assaults, is surely within his good
right when he now in this first section concerning ``the relation
between the individual elementary functions'' expects to find,
not empty accusations and insipid remarks, but clear
proofs, from which it irrefutably follows that here really
``a confusion of concepts'' takes place. But wherein is
the confusion of concepts to consist? where is the proof? That two functions are now separated, now again ``thrown together into one'',
from that surely no confusion of concepts can without further ado be inferred. Multiplication and division, product and quotient,
are, to mention an example, when seen from a
% --- p. 19 ---
\opage{19}\setcounter{footnote}{0}%
viewpoint to be conceived as two separate functions; for, if they
were not two separate ones, how could they then be regarded as
inverse, mutually canceling functions? Nonetheless
these separate functions are yet in another sense to be thrown together into one; for if they did not go together
into one, how could they then be designated by one and the same
fundamental equation? Not thus, but yet in another
manner can the power $x^{a}$ and the exponential function $a^{x}$
be ``thrown together into one'', although they are two independent functions. Insofar, namely, as the original
method of calculation is at issue from which they proceed, they must infallibly be ``thrown
together into one'', for they both arise through the transition of multiplication to potentiation, and insofar both belong
without difference under ``the category of the power'' (``Phil. and Mathem.''
p. 36); but in their fundamental equations the said
functions show themselves, on the other hand, as two independent ones, and are therefore
in ``Phil. and Mathem.'' p. 9 separated by corresponding
fundamental equations. One would not believe that a mathematician, and least of all a mathematician who out of pure
sense of duty has come forward in order to point out ``gross errors'' and
``lamentable confusion of concepts'', and to serve precision, could
hit upon getting stuck on such simple things, and then on top of that run aground on such simple things, and, strictly speaking, himself
be left hanging. ``In mathematics'', says Prof. Steen,
``the transition from ``multiplication'' to the exponential function is not made in the manner indicated; the transition meant here
is made to the power, not to $a^{x}$, but to $x^{a}$, and
not from the concept of \emph{multiplication}, but from the concept
of \emph{product}, for multiplication is an operation of calculation, the product
$ax$ a function''. See, this Prof. Steen now says
decidedly, precisely about mathematics; we must, however, not
forget that it is Prof. Steen who says it, and not
mathematics; for it could after all be possible that here in
what Prof. St. so decidedly and precisely says about mathematics some ``confusion of concepts'' showed itself, and if here
% --- p. 20 ---
\opage{20}\setcounter{footnote}{0}%
now there really were ``a confusion of concepts'', and we then in distraction came to confuse mathematics with Prof. Steen,
and Prof. Steen with mathematics, so that we, instead
of referring them to ``two independent functions'', got them
``thrown together into one'', then it would after all produce an even
greater confusion of concepts.

That there must be ``a confusion of concepts'' in the correction dictated by Prof. St. here,
I conclude from two things: in the
first place from this, that it stands in Ramus, in the second place
from this, that it lies in the nature of the matter.

The elementary theory of the individual functions $a + x$,
$a x$, $a^{x}$, $\log x$ is namely in ``First Section of Algebra and
Theory of Functions'' laid out on the following articulation: ``\textit{§ I.} Sum
(Addition): $a + x$. \textit{§ II.} Product (Multiplication): $a x$.
\textit{§ III.} Power (Raising to a power): $a^{x}$. \textit{§ IV.} The logarithm:
$\log x$, and finally \textit{§ V.} Functions; Fundamental equations: $f(x)$.'' This articulation set up by Ramus shows
clearly enough that ``the transition meant here'' is made precisely from
the concept of multiplication, in that it is made from the concept of product, and that it is made just as much from multiplication to
``the exponential function'', as from the product to the power, in that ``the power'', instead of being designated by $x^{a}$, is precisely
designated by $a^{x}$. If Prof. St. should perhaps regard the
articulation appealed to as something peculiar to Ramus, and in honor of ``the free movement of thought'' appeal to
``the mathematicians of Europe''; then we shall, so as not to be overwhelmed
 by so exalted a notion as the notion of
``the mathematicians of Europe'', again let ``the free movement of thought''
go as one with ``the necessary movement of the subject matter'', and appeal
to sound reason. If in any way there is to be talk
of making a transition whereby the one function is derived
from the other, one must necessarily keep to the peculiar operations of calculation in which the functions originally have
their becoming and coming-to-be; if the concept of function, in a
clarification of this becoming, on the other hand, has already been led to
% --- p. 21 ---
\opage{21}\setcounter{footnote}{0}%
the turning point that the fundamental equations for the four individual
functions $a x$, $x^{a}$, $a^{x}$, $\log x$ can, as Ramus says,
``for each of these functions express the law according to which
the successive values of the function follow upon one another''; then
there can no longer be any talk of a transition, for then
the functions are ``\textit{irreducible}'' to one another.

Thus: if weight is to be laid on the separation of $a^{x}$ and
$x^{a}$, then a transition from ``product'' to ``power'' can be made just as
little as a transition from ``multiplication'' to ``exponential function''. Yet what do I say, that such a
transition cannot be made? Prof. Steen will surely undertake
to make the transition! But is Prof. Steen ---
perhaps with the help of ``Cognition and Investigation'' --- now
also really in a position to make the transition? Yes, if
Prof. can make the transition, can perform the remarkable
feat of making the transition from the concept ``product'' to
the concept ``power'', without it, mark well, --- which
he with his great precision must surely himself regard as a sign
of ``confusion of concepts'', indeed, of ``lamentable'' confusion of concepts, --- at the same time becoming a transition from ``the concept of multiplication'': then I, who ``neither am nor wish to be
a mathematician'', surely do not know what admiration such a
transition will arouse among ``the mathematicians of Europe''; but
this I know, among philosophizers, great and small, the
transition will arouse amazement. A product without multiplication, a multiplication without product! the thinker who
can hold fast such a separation, and make such a transition, is raised above ``confusion of concepts'', has enough in ``Cognition and Investigation'', and can disdain to engage
with ``the dialectic.''

II. ``He (the Author) has not rightly understood the manner
in which the mathematical concept of the imaginary arises, and
there prevails a good deal of obscurity with regard to other forms
under which magnitude appears in ``the Elements''; here is the complaint
grounded in the second section of ``the Critique'': we will
% --- p. 22 ---
\opage{22}\setcounter{footnote}{0}%
still follow Prof. St., in order to examine whether this
complaint might not possibly have some scientific ground or other. That an expression such as $\sqrt{-1}$ or $\sqrt[2m]{-a}$ can
be an impossibility for thought, and that the remark is made (Phil. and Math.
p. 38) that ``an imaginary magnitude
belongs by definition to the impossible'', these, as it
might seem at first glance, rather innocent utterances
have aroused Prof. Steen's indignation even in the highest degree.
``Herein lies'', he says, ``a \emph{complete misunderstanding of the elementary development of concepts in}
mathematics''. This accusation is loud-voiced enough;
but the proof? ``Why'', says Prof., ``should it be
more possible to subtract 4 from 3 than to find the square root
of --- 1?'' We answer: why should it indeed be more impossible? In olden days it was certainly impossible to subtract
4 from 3, because in olden days it was impossible to take 4
where there were only 3, to pay out 4 rix-dollars when one only
had 3 rix-dollars. Has this impossible perhaps now in more recent
times, through the progress of mathematics, become a possibility?
Can the new doctrine of the negative perhaps teach one how
to go about it when one wants to pluck hair from the
bald, or get 4 rix-dollars out of a pocket that only
holds 3 rix-dollars? To find the square root of --- 1, is,
insofar as one thereby expressly understands finding a number
which, raised to the second power, gives --- 1, by many, and
those even very sensible people, by people who otherwise do not
``place themselves on the standpoint of immediate loose consideration'', regarded as impossible, just as impossible as to construct
a round square: can Prof. Steen by his new doctrine
of the negative now really open for us prospects of getting
the round square constructed? if he possesses in his exact
analyses the art of making the in itself impossible possible, then
he ought to say it straight out, he ought to make it publicly
known. He ought to make it known by beat of drum,
% --- p. 23 ---
\opage{23}\setcounter{footnote}{0}%
and for two reasons: in the first place, that we may admire
him as he deserves, for he who makes the impossible
possible deserves admiration; in the second place, that we may at once
concede to him everything, so that we do not go and prostitute
ourselves by unseemly objections, refutations, and all sorts of
intellectual quackeries, because it pleases such a man
to ascribe to us ``a complete misunderstanding of the
elementary development of concepts in mathematics.''

But, someone will perhaps object, here it is surely not really
a matter of making the impossible possible, but of
demonstrating the possibility even there where the unskilled only
all too easily sees the impossible. Might it perhaps be
this in itself very sensible proposition against which ``the Author'' here has
sinned? The reader who has not yet quite
forsworn engaging with ``the dialectic'' will then, as regards
the question of the possible and impossible, not overlook what stands expressly in ``Phil. and Mathem.''
p. 38: ``If with the logical proposition: \textit{possibile}
\textit{est, quod non implicat contradictionem,} one could, as regards concepts, keep to nominal definitions, one would
soon be finished with the separation of the possible and
the impossible; but a thorough separation is conditioned by
systematics\dots{} Yes (p. 39), so dialectical is the limit of possibility, even in the mathematical domain, that
the exact science itself in many points can still scarcely hold it fast.'' Here, then, in ``Phil. and
Mathem.'' the greatest possible concessions have been made to possibility; these concessions are then, as regards in particular the imaginary magnitude, so comprehensive that one must also on
this occasion be amazed at the unreservedness
with which Prof. Steen proclaims ``a complete
misunderstanding of the elementary development of concepts in mathematics.''

``But,'' someone might say, ``the concessions that ``the Author''
thus makes concerning the imaginary's possible meanings
% --- p. 24 ---
\opage{24}\setcounter{footnote}{0}%
do not help him at all; no, it does not help him!''
Why does it not help him? ``Prof. Steen says it
expressly, and what Prof. Steen thus says expressly,
one can rely on unconditionally, Prof. Steen says: it does not help
him, it avails him nothing!'' ``It does not help him
(Contr. p. 8) that the imaginary links the trigonometric
functions to the exponential (Pag. 38), it avails him
nothing that he knows something of the geometric representation of imaginary magnitudes, ``the impossible and for that reason imaginary remains thus still ``in no way real (cf.
the note p. 39); the real and the imaginary must, after
as before the acquisition of such knowledge, be opposite
mathematical concepts.'' The objection consists thus in this,
that ``the Author'' cannot make the imaginary straightway into the
real, and for that reason must still regard the real and the
imaginary as ``opposite concepts'': in truth
a striking objection! He who receives a blow from such
an objection gets a sensation as if he had received a
blow on the ear. ``Such a blow is at once an
imaginary and a real one'', says the dialectic; but, it goes without saying, the imaginary and the real continue none the less to be
opposite concepts, thus: ``it does not help, it avails nothing''.
``Such a blow can bring one to see something black'', says
the dialectic, ``can bring one to see a blackboard, and before the
blackboard a trembling pupil, and beside the blackboard a
gruff teacher, mad in his cubbyhole, writhe, slap himself on the belly, shriek: mad, stark mad! it does not help, it avails nothing!''
this too is an imaginary and yet a real one; but, it goes without saying, the imaginary and the real continue none the less
to be opposite concepts, thus: ``it does not help, it
avails nothing!''

``He who states the imaginary to be impossible by
definition is'', says Prof. (Contr. p. 8), ``ensnared in
a trap he has laid for himself''. If it could now
placate Prof. St., were we to refer him to logic,
% --- p. 25 ---
\opage{25}\setcounter{footnote}{0}%
naturally not to ``the dialectic'', but to ``the correct logic'':
we would, since what is at issue here is a nominal definition,
gladly remind him that a definition requires, besides
\textit{a genus proximum, also a differentia specifica.}
If now ``magnitude'' is \textit{genus proximum,} then it is clear that
the \textit{differentia specifica,} which is designated by the adjective:
real in the expression: real magnitude, solely and alone receives its
meaning through the opposition to the \textit{differentia specifica,}
which is designated by the adjective: \emph{\textit{imaginary}}, so that the expression:
\emph{\textit{imaginary}} \emph{magnitude} originally means a magnitude that
is not and cannot be real. Let us imagine for a moment the expression: ``imaginary'' exchanged for ``impossible'', then not even he whom a hysterical \textit{furor uterinus} might otherwise drive
to extremes on behalf of ``the imaginary'', should surely be offended and fall into
convulsions, because a development of concepts (Phil. Math. p. 38),
just when it aims at possibility, circumspectly inserts the
intermediate proposition: ``but the impossible magnitude belongs by
definition to the impossible: how is this now to be understood?'' If it could now help, were we to furnish proof
that imaginary and impossible magnitude have originally been
synonymous, indeed if it could bring Prof. Steen to
reason, were we, in order to venture the utmost, to appeal to
``the mathematicians of Europe'': then we should, since a lesser thing
cannot do it, gladly name Newton to him. When Newton \textit{(Universal Arithmetic, Part. II. p. 355---56 \&c.)} discusses the equation: $x = a \pm \sqrt{aa - bb},$ arising from
the solution of $xx - 2ax + bb = 0,$ he says: when
$aa < bb$ the roots become ``\emph{\textit{impossible}}'', for the even root of
a negative magnitude is \emph{\textit{impossible.}}\footnote{\textit{\dots{} but when a a is less than b b, they become impossible; because then a a --- b b will be a negative Quantity, and the square}
\textit{Root of a negative Quantity is impossible.}
(Work cited, p. 356).}

But when it is now incontrovertible that the word ``imaginary'' originally has in mathematics reference to the impossible;
% --- p. 26 ---
\opage{26}\setcounter{footnote}{0}%
what semblance of ground has Prof. Steen then for throwing
himself upon so natural and well-motivated a transitional proposition
as this: ``but the imaginary magnitude belongs by
definition to the impossible, how is this to be understood?''

Should anyone still be of the opinion that mathematics in more recent
times has made possible what Newton here declares
impossible, we can let Newton himself enlighten him about the error. As an example that can show how the mathematical language uses ``the imaginary'' to designate
the impossible\footnote{That Prof. St. cannot be ignorant of this (cf. e.g. ``Elements of Analytic Plane Geometry'', p. 53), goes
without saying; but why then does he kick up such a smoke about nothing?}, he refers to the problem of determining
the intersection between a straight line and a circle.

For if we call the straight line's distance from the center
$p$, the radius $r$, the chord that is cut off $c$, then we have
in answer to the question: how great is the chord? the reply:
$c = 2\sqrt{r^{2} - p^{2}}$. If now $p > r$, $c$ becomes imaginary,
or, as Newton expresses it, impossible.\footnote{\textit{In Algebra the Root of an impossible Quantity has}
\textit{its Expression; but in Geometry none. In Algebra you}
\textit{obtain a general Solution, and there is an Expression in all}
\textit{Cases of the Thing required; only within certain Bounds, that}
\textit{Expression represents an imaginary Quantity, or rather is the}
\textit{Symbol of an Operation which in that Case cannot be performed; and serves only to shew the Genesis of the}
\textit{Quantity, and the Limits within which it is possible.}
(Work cited, p. 357, note). Cf. Chasles. \textit{Géométrie supérieure.}
$p. 57$ and especially $p. 547$.} The imaginary $c$ thus signifies the impossibility of the line's being able
to cut off any chord from the circle. But why is
this impossible? because the line lies wholly outside the circle.
But then there is here after all a line, and we can from another
side see something real connected with the imaginary.
For by transforming $2\sqrt{r^{2} - p^{2}}$, for $p > r$, into
$\sqrt{p^{2} - r^{2}}\sqrt{-1}$ one sees that the imaginary magnitude
% --- p. 27 ---
\opage{27}\setcounter{footnote}{0}%
has a real component, and that this part is precisely the
piece of the tangent to the given circle which a circle tangent to the given
line cuts off.

In order now to see quite distinctly how the impossible and
the possible are determined by one another, in that the imaginary
alternates with the real, we will consider the example cited
a little more closely. Suppose the problem is posed in the
following form: There are given two concentric circles, the
one with radius $p$, the other with radius $r$; find
the length of a chord in the outer one which is tangent to the
inner one. In the solution of this problem we must then distinguish
between:
$1$) $p < r$, so that one gets:

\begin{center}$c = 2\sqrt{r^{2} - p^{2}}$ ($I$)\end{center}

$2$) $p > r$, so that one gets:

\begin{center}$c = 2\sqrt{p^{2} - r^{2}}$ ($II$)\end{center}

By multiplication by the imaginary magnitude $\sqrt{-1}$,
one can now easily reduce ($I$) to ($II$).

From this it follows that ($I$) is impossible when ($II$) is possible,
and that conversely ($II$) is impossible when ($I$) is possible. This
mathematical result is moreover in perfect agreement
with the logical, for $p$ cannot possibly, --- in such cases
the dialectic knows very well what it must concede to formal logic, --- be at once both greater (to which
($II$) gives the answer) and at the same time (to which ($I$) gives the answer) less than $r$.

Having given these explanations, I can now join Prof. Steen
insofar as I can conclude with his own words: ``These are familiar things, but one
is compelled to repeat them when one intends to be understood
also by the less initiated''.

III. ``He (the Author) is completely at sea with
regard to a very important question in the integral calculus,
concerning the significance of the arbitrary constant, the manner
% --- p. 28 ---
\opage{28}\setcounter{footnote}{0}%
in which it is determined, and as a consequence thereof also concerning
definite and indefinite integrals''. (Contr. p. 48; treated
pp. 9---19).

Thus: ``completely at sea!'' To point out errors is
a sign of superiority, to point out many errors a sign
of much superiority; we shall here, in order to illuminate Prof.
Steen's superiority, merely bring out some points,
for to go into particulars is surely not worth the trouble.
``I have really'', says Prof. (Contr. p. 12) ``with a
good will sought to meet the Author's explanations''; this
declaration is animating, for it is orienting. For just as
the whole Contribution in general, so is this
section in particular excellently suited to show us what
``a good will'' is capable of. The following examples are characteristic. Against the difference stated (Phil. and Math. p. 76)
between the broken and the curved line ``the
good will'' (Contr. p. 11) puts forward the priceless correction that
the curved line is composed mathematically of finite
lengths just as well as the broken and straight one''. What is adduced (Phil.
and Mathem. p. 79) against the possibility of ``piecing together the triangle out of infinitely many small rectangles''.
is with ``a good will'' rejected (p. 11), as if this
really could be done by virtue of the mathematical concept
of the infinite, according to the infinitesimal method''\footnote{Can Prof. now really show us how one in the infinitesimal method goes about such a piecing together, how
it can happen without one, mark well, (as Phil. and Mathem.
p. 70 expressly emphasizes) transforming the problem by
applying the integration resting on the nature of the function: % print: Functioneus (misprint)
then it will be a performance as brilliant as when he holds fast the function: product after having cast away the concept: multiplication. But if he cannot do this, he must himself admit that one
according to the infinitesimal method just as well, when it comes down to it, must make use of the relation between a derived and an original
% (the note runs on to p. 29)
function, as in what he calls ``the more stringent presentation'': what right has he then to put forward his correction, what
ground has he for his ``mathematically incorrect''? The remark, in itself as incontrovertible as it is easily grasped (Math. Phil.
p. 67), that one cannot without further ado set $dx = 0$, or
$\int dx = \infty 0$, because $\infty 0 = \frac{1}{0} 0 = \frac{0}{0}$ is, both as integral and
as differential coefficient, wholly indeterminate gets ``the good will''
by a couple of curious twists finally wrested into the following bungling: ``the Author'' draws the conclusion that when two expressions
occur under the indeterminate form $\frac{0}{0}$ (of which the one is yet
originally $\infty 0$ and only by a formal application of
the sign of equality becomes $\frac{0}{0}$), then there is no difference
between them!'' To this the answer is short and sweet: if Prof.
can get something determinate out of forms such as $\infty 0$ and $\frac{0}{0}$, thus hold fast a difference between them, without looking to their origin from some function or other: then he performs his third feat, as brilliant as the separation of multiplication and
product, and as surprising as the piecing together of the triangle
out of infinitely small rectangles, and we shall admire him; but if
he cannot do this, we shall not admire him, for then his correction
is merely a schoolmastering, and such a schoolmastering wholly unwarranted.
That the forms $\infty 0$ and $\frac{0}{0}$ cannot be distinguished or more closely
determined when one abstracts from the function, is just the thought
which in Phil. and Mathem. (p. 67) is expressly emphasized:
% (the note runs on to p. 30)
``Mathematics does not engage in determining $dx$ immediately, be it as point or line; $dx$ is determined only in a
relation as e.g. of $dy$, so that the law of change'' --- namely
for $dy$ in its dependence on $dx$ --- ``can show itself in $f(x) = y$.
In ``Contr.'' p. 13, ``the good will'' advances with the following
sally, which for that matter might just as well be called
a whim. ``But let us now'', says Prof., ``assume that
it is correct that $\int dx$ cannot be regarded as a determination of $\infty \cdot 0$'' (in the form of a concession the remark
is false, for in Phil. and Mathem. $\int dx$ is expressly regarded
as a determination of $\infty \cdot 0$), ``and that the area $\int y dx$ cannot be
pieced together out of infinitely many infinitely small rectangles, what
then is ``the infinite'' that is brought back by the integration to
``finite determinacy''? Then there is in $\int y dx$ nothing mathematically
infinite and then the infinite, if it has any significance, must have a philosophical significance --- truly a muddle
in which one had best refrain from trying to find a clue!'' Thus: because one
cannot ``piece together'', mark well, ``piece together'', says and
writes: ``piece together'' the finite out of the infinite, there is in
$\int y dx$ nothing mathematically infinite: truly, a muddle as
complete as anyone could wish for, who merely wishes for a muddle,
because he merely wishes to point out errors; yes, to point out errors is a
sign of superiority, to point out many errors a sign of
much superiority, to point out tremendously many, tremendously great and tremendously
gross errors a sign of a tremendous superiority!}

% --- p. 29 ---
\opage{29}\setcounter{footnote}{0}%
When it is now moreover proclaimed in conclusion, in lapidary style
and with due trumpet blasts, to all, to the initiated as well as to the ``less initiated'' and to the uninitiated, that
the ``Author'' in Phil. and Math. is ``completely at sea
with regard to a very important question in the integral calculus, concerning the significance of the arbitrary constant,''
then the initiated naturally grasp easily what such an announcement is meant to signify; but the less initiated? Well,
the less initiated, in particular ``the young students,'' might
perhaps, by the strong trumpet blasts, especially when they hear
that the announcement even proceeds from ``the good will,''
easily get at sea concerning ``the arbitrary constant,''
and all the more easily as ``the good will'' is, after all, also in
% --- p. 30 ---
\opage{30}\setcounter{footnote}{0}%
certain cases ``an arbitrary constant.'' Let me then, to
be on the safe side, point out to ``the less initiated'' that in the integral calculus there is no difficulty at all with the ``arbitrary constant'' that has been invoked.
For if we denote by $dx$ an infinitely small increment of $x$,
then it follows of itself that $x$, at the moment when it receives
the increment, must already have a certain determinate value, be
it $0$ or $a$ or $c$, and thus be conceived as a constant.
From this simple presupposition consequences can then be developed in the most various directions; here there are not merely
one but a hundred ways in which the significance of ``the
% --- p. 31 ---
\opage{31}\setcounter{footnote}{0}%
arbitrary constant'' can be presented and explained. If now
Prof. St. nevertheless insists that there is only
one single correct way, namely his own way\footnote{By setting an area like $bdBD = \int_{x^{1}}^{x} y dx = f(x) - f(x^{1}) =$
$f(x) + C$, ``every distinction
between definite and indefinite integral is entirely effaced,'' according to Prof. St.'s assertion; but this
is of course merely an assertion, since, unless an arbitrary dictum forbids all use of incomplete equations, one
can interpret the formula given thus: the area \textit{bdBD} now shows
what it means that $\int_{x^{1}}^{x} y dx = f(x) - f(x^{1})$ is a closer
determination of $f(x) + C$, in that $C = - f(x^{1})$, in other words, that
the definite integral is precisely a special case of the indefinite.}: then
let him announce it with trumpet blasts, and he will find
mathematicians, even mathematicians of the first rank, in
``{}``completely at sea concerning the significance of the arbitrary constant,
the way in which it is determined, and as a consequence
thereof also concerning indefinite and definite integrals.''

``I regret,'' says Prof. Steen as he concludes
his 3rd section, ``having had to weary the reader for almost
four pages in order to set forth on one page what is erroneous, meaningless and
unclear.'' What exertions the Professor's
``good will'' has had to make in order thus to ``weary
the reader'' can be seen clearly enough from the means he has had to
employ. To enable even ``the less initiated'' to
appraise the nature of these means, a single
example will suffice. In order to make intelligible why
one must always, when, as Prof. Jürgensen puts it,
one ``wants to give the integral the most complete form,'' add
a constant, even in such cases where the constant,
when more closely determined, proves to be $0$, there is
in Phil. and Math. (p. 82) a direct reference to the determination
of the area of the triangle and thereby indirectly to the determination of
the trapezium.

% --- p. 32 ---
\opage{32}\setcounter{footnote}{0}%
Without the slightest regard to the determinate train of thought of the treatise\footnote{For before the example with the calculation of area there are given (p. 81)
a couple of other examples, from which it is unmistakably seen that ``the Author''
is not at all at sea with the arbitrary constant and the
way in which it is determined; but these examples Prof.
Steen, ``with a good will,'' has of course passed over in complete silence.} Prof. St., with a dogmatic insistence on being right that
can surely without exaggeration be called violent, introduces his own train of thought, which has nothing to do with the
real question, and thereupon begins with the following twistings, as he so aptly
puts it, ``to weary the reader.'' Is the area of the triangle determined by $y \cdot \frac{x}{2}$ from the point of view of finitude,''
it is said (Contr. p. 15), ``then from the standpoint of infinity one must apply the formula: $\int y dx + C$.''

Now the corrections follow. First correction: ``It is a
lack of precision when it is not stated here what $x$ and
$y$ signify, and one can only guess that the remark
refers to Example 2, p. 79.'' It is ``a lack of precision'' in this
correction that it contains
a plain untruth, for the figure to which the exposition
refers is found (p. 82) in immediate connection with the
explanation itself; the correction has, next, as regards the concept of precision, the defect that, even if it were
entirely precise, it is pedantic. Such corrections are no doubt in
the highest degree unfortunate when one seriously aims to
``promote the dissemination of scientific truths,'' but
fortunate enough when the task is to garble a scientific
development, and to ``weary the reader.''

Second correction: ``If $\frac{yx}{2}$ is taken to be the area of the triangle, then $\int y dx = \frac{xy}{2} + C$ is only in one case
the area of a triangle, namely when $C = 0$, but then it would
% --- p. 33 ---
\opage{33}\setcounter{footnote}{0}%
have to be stated that $\frac{yx}{2}$ was the area of the triangle
from the point of view of finitude, $\frac{xy}{2}$ on the other hand from that of
infinity! If one chooses this conception, then one
gets the same result from the two different
points of view; but to what end, then, a juxtaposition of them,
as if there were a difference?'' This correction is unforgettable, for it is \textit{colossal}\footnote{``No one,'' says Prof. St. (Contr. p. 15), ``who has followed the
preceding development with attention, will be able to deny that
on the three leaves here discussed in the first instance, pp. 75---80, there prevails
obscurity and confusion with no small admixture of incorrectness; but that is nothing compared with page 82, where there swarms with errors
that fully demonstrate that the Author has either not understood at all
the simplest propositions of the integral calculus or does not have the mastery over them that ought to prevent him, when he is to give
independent developments, from committing the very grossest errors.'' Yes, ``the
very grossest errors''! Beaten by such ``very grossest errors!'' ``the Author'' must
console himself with a little piece of moral from Claudius, in which it says so
consolingly: ``Ein großes Maul es auch nicht thut.''}. For in ``Phil. and Math.''
the fundamental thought is this: the same magnitude, the
same area, e.g. that of the triangle, can be conceived from two
mutually opposed points of view, that of finitude and that of infinity. From the point of view of finitude the area of the triangle is
given by the formula $y\frac{x}{2}$, since, as the figure shows, $x = OD$
is the base, $y = BD$ the height. Here there can be no question
of adding any constant, and why not? Because
$x = OD$ and $y = BD$ are finite, arbitrary magnitudes, which are
chosen as existing and conceived as given. From the standpoint of infinity, on the other hand, the determination of the area of the triangle must necessarily fall under the general
formula: $\int y dx = \frac{xy}{2} + C$, that is to say: here emphasis must be placed
on the addition of a constant, notwithstanding that the value of
this constant is $0$ in the case of the triangle. The point
% --- p. 34 ---
\opage{34}\setcounter{footnote}{0}%
is that the constant, seen from this point of view, has an equally % print: Synspuukt (misprint)
essential validity when it is $O$ as when it is not
$O$. For if we conceive the magnitudes $y$ and $x$ as infinitely
variable in such a way that $x$ in $f(x) = y$ is the independent, $y$
the dependent: then the area is conceived, according to the formula:
$\int y\,dx = \frac{xy}{2} + C$, in such a way that its existence is expressly
seen as the result of a coming-to-be.

But what has Prof. Steen's ``good will'' not now
managed to get out of this in itself unremarkable exposition, what conclusions has he not drawn, what corrections has he not produced! Imagine an assembly of ``the mathematicians of Europe,'' imagine
Prof. St. appearing in the assembly, delivering his corrections,
and each time hitting the point: what tension of
attention, what listening, and then --- what surprise!
``If one gets,'' so runs Prof. St.'s correction, ``the
same result from two different points of view, to what end then
a juxtaposition of them, as if here there were a difference?''
that is to say: when two different roads lead to the same goal,
then there is no difference between the two --- different roads:
a fine remark, a lovely point! ``If $C$,'' the Professor continues, ``is not $O$, then $\int y\,dx = \frac{xy}{2} + C$
is the area of a trapezium, and then it would have to be stated that
$y\frac{x}{2}$ is the area of the triangle from the point of view of finitude,
but from the standpoint of infinity one must
for the area of the triangle apply an expression for
the area of a trapezium!'' Thus: since
$\int y\,dx = \frac{xy}{2} + C$ is not a triangle if it is a
trapezium, and not a trapezium if it is a
triangle, these two cases are to be conceived in such a way that
the latter case is merely a special case of the former
case, for when $C$ is $O$, then this is (Contr. p. 16)
of course merely a special case of the case that $C$
% --- p. 35 ---
\opage{35}\setcounter{footnote}{0}%
is not $O$: how acute! ``Thus,'' Prof.
St. continues, ``the philosophical consideration here (Phil. Math.
p. 82) has brought us into conflict with the logical proposition: \textit{possibile}
\textit{est, quod non implicat contradictionem}, according to which the Author
nevertheless thinks on p. 38 that a round square is impossible, but on p. 82
arrives at the result that a quadrilateral area from the point of view of infinity
is a triangular one.'' But --- let us hear it once more ---
how has ``the Author'' actually arrived at the result that a quadrilateral
area is triangular? ``Why, that is clear from what Prof.
St. says, and Prof. St. is always right in what he
says in such a way; he has of course arrived at it by
assuming, \textit{plaudite jam}, in his gross ignorance, $C = O$
to be a proper constant.'' --- A pity that $C = O$
is not a proper constant; for, it goes without saying, if
zero, pure zero, were a genuine constant, then the triangle
would surely refrain from being quadrilateral, and if $C = O$ could first
keep the triangle from being quadrilateral, then $C > O$
would in return soon cure the quadrilateral of its whims of
wanting to be triangular! Lo, with this sort of entertaining
inferences a Prof. St. argues; and now let
some among ``the mathematicians of Europe'' step forward and say that
a Prof. Steen is not as much at home in his logic as
in his constants; now let some say that a Prof. Steen
has not had the benefit of ``Cognition and Scrutiny''; now let
some say that a Prof. Steen with his precision,
his rare, almost incredible precision, must ``weary the reader''!
for that he himself regrets having to weary the reader is surely
``in all modesty'' an allusion to the fact that here, with
his ``good will,'' he --- ``permits'' himself to entertain the reader.

IV. ``He (the Author) has grasped the first mathematical
significance of the concepts zero and infinite, as pure limits of magnitude
. On the other hand he has only an unclear notion
of the essential character these concepts have as
means of expressing negations in positive form, and finally he has no inkling at all of their significance as
% --- p. 36 ---
\opage{36}\setcounter{footnote}{0}%
transitional forms, from which there shines forth precisely something qualitative
in opposition to pure quantity, within which mathematics otherwise moves. But herein mathematics and philosophy have precisely their most essential point of contact,
concerning which it surely holds that it could be fully illuminated
if the two sciences were to meet each other halfway \dots{}. He
is partly, as a consequence of an incorrect and unclear conception
of the concepts zero and infinite, at sea with regard to
all formulas in which these peculiar forms appear,
such as in giving $\frac{0}{0}$, $\infty 0$, $\frac{\infty}{\infty}$ etc.'' (Contr. pp. 48---
49; cf. pp. 19---30).

The mathematical creed with which Prof.
St., as has already been remarked, introduces this section,
in that he bears witness to his humble opinion
concerning the problem of infinity, presents on closer
consideration several rather interesting aspects. Just as it
is a categorical error to want to introduce an objective
knowing where the subjective assurance of faith must be the
only thing valid, so it is on the other hand an equally
categorical error to want to apply faith where an
objective knowing alone can be valid. That Prof. St.
with his ``rigorous critique'' really has an inclination, and one that is moreover
especially very questionable in a mathematician, to introduce, where
it suits him, the comfortable categories of faith, of opinion and of the
subjective conception, is clear from many indications. When he thus, on the occasion of the Ramus --- and, alas,
also Lagrange --- ``quite peculiar proof'' (Contr. p. 40)
utters, with a modesty by no means unbecoming to him but certainly unusual:
``I permit myself to add some doubt whether
the mathematicians of Europe will be able to agree to acknowledge
Ramus's proof as completely satisfactory?'' then one gets
the impression that the mathematical theorem had been
transformed at a stroke into an article of faith, as if the validity
of a proof could be decided in the highest instance for the individual
% --- p. 37 ---
\opage{37}\setcounter{footnote}{0}%
mathematician by his having it sanctioned at a council of mathematical notabilities, just as the believing
Catholic has his dogma sanctioned at a council of
ecclesiastical authorities. Here, then, is after all a faith in authority
that involuntarily stirs in Prof. Steen, and that notwithstanding that it is precisely Prof. Steen who on this occasion has felt himself called upon to appear ``not with any
self-assumed or imagined authority, not with reference to
other mathematical authors'' (of course with the exception
of ``the mathematicians of Europe''), ``but solely with
reasons.''

In Prof. Steen's conception even of the merely mathematical one notices not only the specialist, who one-sidedly
sees everything in the light of his own particular science, but at the same time the mathematical man of business, who has already acquired a
definite manner, a somewhat stereotyped procedure, in which
routine is to make up for lack of freedom of mind, indeed
the complete man of authority, who demands respect for his opinion, his individual opinion, as one demands respect for
a trump. If one may not proceed from such psychological
presuppositions, then the striking unabashedness, the
extraordinary self-confidence with which a man who after all
does not wish to appear ``with any self-assumed or imagined authority,'' in the name of mathematics dares to brand every proposition that
does not please him with his ``incorrect,'' ``wrong,'' ``stark wrong,''
``meaningless,'' as if he were not merely a mathematical specialist
but mathematics incarnate in his own person, becomes altogether
inconceivable; only with a strong admixture of something belonging to faith in the specialty
does it become explicable that a man whose education
has hardly led him in the direction that most favors
a stricter, more many-sided and coherent development of concepts,
blithely sails in against what he demonstrably has not understood,
indeed, as it seems, has not even taken the slightest trouble
to understand, but rather --- lest we fail to appreciate
``the good will'' --- the trouble to misunderstand.

% --- p. 38 ---
\opage{38}\setcounter{footnote}{0}%
That Prof. St., with all respect for his routine
as a master calculator, is by no means the man to turn to when one wants information about the infinite in mathematics is clear from his whole Contribution, and becomes
especially evident from what he sets before us in this and in parts of the following sections. Already the beginning he
makes by conceding that there is indeed a mathematical
concept: infinite, ``which arises and comes to be applied,
and which is extremely important and moreover difficult, because it is
of a mixed nature,'' lets us surmise that in the infinite, which he
regards as extremely important and moreover difficult, we shall meet with
a good deal of hodgepodge that one had best see to be rid of, so as
not to encumber the development of concepts. Right at the moment when he lets his
schoolmasterish critique swing the corporal's cane over the first
series drawn up, he strikes wide of the mark. In the long correction (Contr.
p. 21) there is not a word, not a single word, that concerns
what he wants to correct. In his zeal and haste he has namely seen
wrongly, and at once loses his way among ``the different roads''; for by
``the different roads'' is by no means denoted what he, who seldom troubles himself about any connection, guesses at, but,
as plainly emerges from the context, $f_{1}(n) = a + nx$
for the arithmetic, $f_{2}(n) = \frac{1}{a + nx}$ for the harmonic and
$f_{3}(n) = ax^{n}$ for the geometric series. To be sure, he has
as it were a feeling that here on the whole there is something odd, and in particular a little ``but'' in the text has permitted itself
to make the Professor aware that he must surely have
gone astray; but Prof. St., who, as said, has
an inclination to faith in authority, and therefore can never believe that it
should be he, he himself, who erred, but on the contrary very
easily believes that ``the Author'' has committed a blunder, were it even % print: eud (misprint)
an enormous blunder, blithely puts his own ``while'' in place of the
text's ``but,'' and so the correction goes off at a run.

``I shall,'' says the Professor (Contr. p. 22), ``leave entirely out of consideration that the connection between all
% --- p. 39 ---
\opage{39}\setcounter{footnote}{0}%
this'' (Philos. and Mathem. pp. 20---26) ``is too loose to
deserve the name of a ``development of concepts''; but the incorrectnesses occurring therein
must be pointed out.'' Thus: the Professor
will no doubt assert that the connection, and that a connection of what occupies fully 6 pages, is too loose to
deserve the name of ``development of concepts''; but if anyone
should happen to doubt his assertion, and demand of him
a little justification, then it is declared at the same time that here there is no reason to ask for any
reason, since such a thing is to be ``left entirely out of consideration'';
thus: two declarations in one declaration! If here
faith in the specialty, faith in authority, and that a rather
blind faith in authority, is not presupposed, what is authority? what is faith in authority? When Prof. St. furthermore (Contr. p. 24), on
the occasion of his opinion about something of which he has not
grasped the meaning, exclaims: ``But should that be the meaning, then it and similar remarks would be doomed to death,'' does he not speak a language as if he sat on
the papal throne of mathematics, proclaimed the dogmas, and hurled his anathema
against ``the doomed''? Yet let us
not forget it, Prof. St. wants no faith in authority,
Prof. St. is precisely the man who has ``felt himself called upon to appear, not with any self-assumed or imagined authority, but solely with reasons.''

An example of what sort of hodgepodge Prof.
St. wants to foist upon philosophy when it ventures to
touch the mathematical infinite is ``the essential character the concepts zero and infinite,'' according to his opinion,
``are supposed to have as means of expressing negations in positive form.''
``When it is asked (Contr. p. 27), how much must be added
to $a$ to give $a$, the answer in words is: nothing is to be added,
thus \emph{negatively}, but the answer in signs is:
$a-a = 0$, zero is to be added, thus affirmatively,
in that the negation is put into the zero. When it is further
asked, how many times must $0$ be made greater to give
% --- p. 40 ---
\opage{40}\setcounter{footnote}{0}%
$a$, the answer in words is: that cannot be produced by
multiplying $0$, but in signs: $\frac{a}{0} = \infty$, $0$ must be made infinitely many times greater to give $a$; the negation lies
here in infinite.'' After communicating this interesting
information the Professor adds triumphantly: ``And on
this particularly essential point the Author, who
wants to write about philosophy and mathematics, has only a
single remark, and that has moreover the whole character of accident.''

Although we gladly concede that these, as the Professor
himself calls them, ``simple and easily grasped explanations''\footnote{In his first example of negation in positive form, namely $a-a$
= 0, the Professor seems especially to be very taken; he exclaims
in the note: ``Compare this simple and easily grasped explanation with the Author's remarks on $a-a = 0$ p. 53. I do not permit myself to understand these''\dots{} ``Compare!'' This
triumphant exhortation must, if I may here introduce a
position in negative form, of course not be taken as an outburst of % print: af af (doubled word)
childish joy at having produced oneself in logic, and thereby produced a simple and easily grasped triviality, but much
rather as a hint toward further instruction. Besides the negation
in positive form there is, after all, communicated as a bonus, in the negation so characteristic from a logical point of view: ``I do not permit myself to understand''
a negation in a quite peculiar original form, namely a negation
in the form of filler. If it is now perhaps at the same time the Professor's opinion that such negations in the form of filler are likewise to constitute essential points in his doctrine of zero and infinity, then
I shall raise no objection to it; but they will hardly benefit
either philosophy or ``the young students.''}
bear witness to deeper ``scrutiny'' in formal logic: yet we can neither attribute to them significance in the doctrine of the
infinite, nor assume that they would just in any appreciable degree
appeal to ``the young students.'' As to the impossibility of
making an infinite distance finite, a single remark,
even if it has the whole character of accident, more than
% --- p. 41 ---
\opage{41}\setcounter{footnote}{0}%
sufficient. If one wants to go further and investigate how
the infinite in different cases comes to denote
now the possible, now the impossible: then it will soon be shown
that the particular significance rests on an admixture of
something that does not in and for itself concern the concept of infinity. If we seek, e.g., the geometric locus of the centers
of spheres that touch three given spheres on the same side, we obtain
a hyperbola; the infinite radius, which here denotes the negation, will then belong to the sphere corresponding to the point at infinity on the
hyperbola. But a sphere with an infinite radius is no sphere, it is a plane: the impossibility of the sphere is to that extent the possibility of the plane. $\infty$ can
here further denote the transition from the positive to the negative,
from the spheres with positive radius to those with negative radius, from spheres in the series that touch the three given ones externally to those that touch them internally.

Prof. Steen's thoughts are unruly: his reflection has incredible whims. Where impossibility, as
in the imaginary $\sqrt{-1}$, is expressly laid in the notation
itself, the Professor will not acknowledge it, will not hear a word about
it, will not have anything at all to say of it; where impossibility,
on the other hand, as in $0$ and $\infty$, lies neither in the concept
nor in the notation, but presupposes such peculiar considerations of
peculiar conditions that it could just as well be denoted by
a positive as by a negative magnitude, just as well by a finite
as by an infinite one, there Prof. St.'s capricious reflection absolutely wants to have the impossible pressed into $0$ and $\infty$;
only in this way can the Professor glorify his discovery, his great discovery, the discovery of ``negations in positive form!''

Since my conception of the relation between the finite
and the infinite has already received a further development in
another work,\footnote{Mathematics and Dialectic. 1859. Hegel's Publishing House.} there shall here in this context only
% --- p. 42 ---
\opage{42}\setcounter{footnote}{0}%
be given a couple of examples, in order to show further
with what self-confidence Prof. St. engages in
judging what he has not at all given himself time to
think through.\footnote{Some examples. 1) Contribution p. 43. Regarding the inference:
``but when it is demonstrable that the ratio for the infinite becomes the same as for the finite, the unattainable can be said
to have been attained,'' it is added: ``This is incorrect,'' and it is classed
even among the examples of ``gross incorrectnesses.'' Such a
correction betrays thoughtlessness; for if the lawfully determined ratio did not remain the same, % print: dct (misprint)
whether the terms of the ratio are otherwise
finite or infinite, the so-called indeterminate forms could never be
determined. The thoughtlessness emerges still more strongly,
as the correction continues: ``already in the first elements of mathematics innumerable examples are met with of one's knowing the ratio of two magnitudes without knowing the magnitudes themselves.'' It takes
only a slight degree of mathematical thoughtfulness to see that
no magnitude can be known without relations and ratios.
What indeed is the unit? If 1 is not to be taken in the same immediate way as one takes a stick: then 1 is to be determined only
as the ratio of two equal magnitudes: $\frac{a}{a} = \frac{x}{x} = \lim. \frac{x}{x} = \frac{0}{0}$
$= \mathrm{Lim.} \frac{x}{x} = \frac{\infty}{\infty} = 1.$ If it now stands firm that every magnitude
is determined by the determinate unit, then it also stands firm that
the unit itself is originally determined as a ratio: the original determination of magnitude is a determination of relation and ratio. 2) In Contribution p. 45, regarding the proposition that the differential coefficients of different order, ``although % print: Differeutialcoefficienterne (misprint)
different from one another,'' yet in comparison with
finite magnitudes, each severally are as $\frac{0}{0}$,'' it is remarked: ``This is
incorrect, for the different differential coefficients are as a rule
\emph{determinate} finite magnitudes''\dots{} ``The error is that the Author p. 55
and 60 sets up $\frac{dy}{dx} = \frac{0}{0}$\dots{}, but forgets that when the differential $dx$ is given, then the differential $dy$ is also given by $dx$.''
% (the note runs on to p. 43)
A single such correction suffices to brand the arbitrariness with which Prof. St. introduces his corrections. In
Phil. and Mathem. p. 55 it is expressly said: ``In relation to any
finite magnitude whatever, $dx$ and $dy$ are both $= 0$, but
to one another they are nevertheless in a quite determinate ratio: $\frac{dy}{dx} = \frac{0}{0}$ just as determinate as e.g. $\frac{y}{x}$ or as $\frac{1}{2}$ or
as $\frac{2}{3}$. That the differential coefficients in this sense can
be \emph{determinate} finite magnitudes, and as a rule are so, is therefore
by no means overlooked; but the point is that the Professor cannot understand what is not set forth in the same words in which he
with his routine is accustomed to get it set forth, and therefore does not
see that $\frac{dy}{dx} = \frac{0}{0}$ here is expressly meant to denote that the determinate
magnitude is one with the determinate ratio, and that this determinate
ratio can be held fast, notwithstanding that the terms of the ratio are vanishing,
that, in other words, --- what he has likewise not understood ---
`the ratio for the infinite, $\frac{dy}{dx}$, becomes the same as for the
finite.' 3) Contr. p. 47: ``Page 75 in the note is found: `Since
a constant cannot have a differential, neither can it
have an integral.' `The first,' it says, `contains no
justification of the latter, and the latter is moreover incorrect.''' --- That
one can now in the machinery of calculations apply the notation for
infinite summation also to terms that are not differentials,
shall be conceded; but in the concept, and here we are indeed dealing with ``development of concepts,'' the integral is so opposed to the differential that
they correspond to one another: integration essentially presupposes differentiation. When it is further said that in $\int a\,dx = ax + C$
$a\,dx$ is constant: then this correction rests on a subtlety;
for by a constant that cannot have a differential, one does
not here of course think of the differential itself: $dx$ can, notwithstanding that it
is constant in a ratio like $\frac{dy}{dx}$, by no means be lumped together with
an immediate % print: umiddelber (misprint)
constant $C$. To go into detail on everything that
% (the note runs on to p. 44)
Prof. St. in unbridled arbitrariness has produced of one thing and another
would be too lengthy. As characteristic of the manifold, entirely unwarranted filler corrections it shall be added only that $\int a\,dx = ax + C$ now nevertheless (Contr. p. 47) has been granted
permission to count as a definite integral, whereas we, on the other hand, Contr.
p. 14 are ``corrected'' with the correction that it is precisely the arbitrary
constant that makes the integral indefinite.}

``It will be most correct,'' he says in Contr. p. 23,
``first to show how the Author's own example the hyperbola and
% --- p. 43 ---
\opage{43}\setcounter{footnote}{0}%
its asymptote, which in illuminating the concept ``infinite approximation'' has nothing at all to do with it\dots{} that the finite
% --- p. 44 ---
\opage{44}\setcounter{footnote}{0}%
should be able to lie infinitely near to the infinite\dots{} for
either the abscissa $x$ is finite, and then $y_{1}$ and $y$ are finite
magnitudes with a finite difference, or $x$ is infinite,
and then $y_{1}$ and $y$ become infinite with the difference zero.'' From
this correction it is seen that Prof. St. not merely for mathematical use, which may be highly commendable, applies
the method of limits, but that he is at the same time, through the great
routine, so bewitched by this method's abstract
\emph{either --- or,} that the method has become a mannerism for him.
But when the mechanical \emph{either---or} of the method of limits has once
become a mannerism for a mathematician, he will surely refrain
from grasping that on which it essentially depends in a
development of concepts of the infinite.

According to the either---or of the method of limits, either finite
or infinite, one comes by a critical leap from the finite to the infinite, just as by a critical leap one can
come from land into water: now $x$ is finite, now $x$ is
infinite,\footnote{Cf. Mathem. and Dialectic. pp. 59---60} now one is on land, now one is in the water.
If one can perhaps by incessant leaping finally manage
to leap over motion in the leap: then one performs
one's fourth feat, an accomplishment as brilliant
as the removal of the concept of multiplication from the function product, as brilliant as the patching together of the
triangle by the infinitely many rectangles, still more
brilliant than the demonstration of a difference between $\infty\ 0$
% --- p. 45 ---
\opage{45}\setcounter{footnote}{0}%
and $\frac{0}{0}$ without application of the law of function; but if one cannot
by any leap whatever, be it logical contortions
or mathematical buck-leaps, leap over motion,
because by no leap can one leap over the leap,
and yet notwithstanding, from sheer sense of duty, feels called upon to
instruct him who sees motion in the leap: then one is
not an authorized teacher, but a thoughtless leaper.

A proposition like that (in Phil. and Mathem. p. 24):
``The distance from the finite to $0$ and from the finite to $\infty$ is
thus here like a distance from the one point of finitude to
the other'' belongs especially among those that must give offense to thoughtless
leapers. As proof of what
enormous error is committed by assuming that an ever so small
motion to the one side or the other brings the tangent of the angle from being infinite to being finite, there is set up (Contr. p. 25) the following equation:
\[ \mathrm{tg}\,90^{\circ} - \mathrm{tg}\,(90^{\circ} \pm \varepsilon) = \frac{\sin (\mp \varepsilon)}{\cos 90^{\circ}\,\cos (90^{\circ}\pm\varepsilon)} = \frac{1}{\cos 90^{\circ}} \]
$= \infty$, ``whereby,'' as it is expressed, ``\emph{the mathematical}
\emph{error is shown to be infinitely large.}'' Without
dwelling on the many fine consequences that
can be derived from the fact that one thus without further ado measures a
mathematical error by the difference overlooked in the error, we shall, in order to test whether there is in fact any
error here, at once ask what it is then that the equation $\mathrm{tg}\,90^{\circ}$
--- $\mathrm{tg}\,(90^{\circ} \pm \varepsilon) = \infty$ is actually supposed to prove. In accordance with the either---or of the method of leaps, either $90^{\circ} \pm \varepsilon$
and then a finite tangent, or $90^{\circ}$ and then an infinite
tangent, an ever so small motion to the one or the
other side will after all bring the tangent of the angle from being
infinite to being finite; and yet Prof. St.
sees precisely in this ``the great misunderstanding'' that he wants to correct by means of the equation set up. Insofar as the equation proves
% --- p. 46 ---
\opage{46}\setcounter{footnote}{0}%
something, it proves the very opposite of what Prof. St.
wants to prove.

That the finite lies infinitely near the infinite will,
if we hold fast to the method of limits, in other words
say that for the finite there is an outermost limit that
cannot be exceeded without the infinite suddenly showing
itself. Such an outermost limit is now precisely the point
in which secant and tangent intersect one another, immediately before
the intersection ceases, i.e. immediately before the angle's
transition from $90^{\circ} \pm \varepsilon$ to $90^{\circ}$, $\varepsilon$ being thought of as the
least possible magnitude. How great is now the distance between the point in which the tangent touches the circle, and the outermost point of intersection of
tangent and secant? \emph{Infinite!}
But then secant and tangent would have to be parallel, and
consequently could not intersect one another. \emph{Finite!} But if
the point of intersection does not lie immediately on the border
of the infinite, it cannot be the outermost; but
if it is the point with which the infinite begins, it must
lie infinitely far away. \emph{Neither finite}
\emph{nor infinite!} But if there really is an
interval between the point with which the finite
ends and the point from which the infinite begins:
perhaps the Leapers could then give us a little information about how the infinitely many tangents
whose points of intersection with the secant fall precisely in this
interval are to be more closely determined\footnote{The increments of the tangent are related inversely as the square
of the cosines to which the angles correspond. If $y = \mathrm{tg}\,x$, then
$dy = \frac{dx}{\cos^{2} x}$. If we thus imagine the one leg of an angle to move from 0 to $90^{\circ}$ with constant angular velocity, then the
point in which the moving leg intersects a line perpendicular to the fixed leg
will move with a continuously increasing
speed (at $45^{\circ}$ this speed is about 2 times as great
as at $1^{\circ}$, at $89^{\circ}$ about 1642 times as great as at $45^{\circ}$);
% (the note runs on to p. 47)
but the transition to these different speeds does not take place by
leaps, but, mark well, continuously. At the transition from
$90 - h$ to $90^{\circ}$, where $h$ is a least possible magnitude, it seems
indeed as if the speed suddenly leapt from finite to $\infty$; but
this lies in the fact that $h$ always indicates a finite discreteness,
insofar as it is not conceived in its vanishing, i.e. in a continued
motion toward $0$. It is with respect to this that in the method of limits we obtain different $\infty$ for $\mathrm{tg}\,90^{\circ}$, according to the manner
in which we let its difference from the preceding vanish.
Consider merely:
$\lim \frac{dy_{90^{\circ}-h}}{dy_{90^{\circ}-nh}} = \lim \left(\frac{\cos (90^{\circ}-nh)}{\cos (90^{\circ}-h)}\right)^{2} = \lim \left(\frac{\sin nh}{\sin h}\right)^{2} = n^{2},$
where
$y_{90^{\circ}-h} = \mathrm{tg}\,(90^{\circ}-h),\ y_{90^{\circ}-nh} = \mathrm{tg}\,(90^{\circ}-nh).$
The only case where the ratio between the two increments
would, strictly speaking, become infinite, would be for $n = \infty$;
but in that case the moving leg would not have moved
continuously from the position determined by $90^{\circ} - nh$ to the one determined by $90^{\circ} - h$;
the motion would then have had to take place by a leap.
The method of limits, which always posits $h$ either as finite or as
$0$, cannot possibly make intelligible the continuity holding in becoming and motion. Posit:
$\mathrm{tg}\,90^{\circ} = \lim \mathrm{tg}\,(90^{\circ} - h) = \infty_{1}$
$\mathrm{tg}\,90^{\circ} = \lim \mathrm{tg}\,(90^{\circ} - 2h) = \infty_{2} = 2^{2}\,\infty.$
$\mathrm{tg}\,90^{\circ} = \lim \mathrm{tg}\,(90^{\circ} - 3h) = \infty_{3} = 3^{2}\,\infty.$
$\dots$
$\mathrm{tg}\,90^{\circ} = \lim (90^{\circ} - nh) = \infty_{n} = n^{2}\,\infty$
Thus $\infty_{n-1} - \infty_{n} = \infty.$
The imperfection of the method of limits shows itself further in this,
that without being able to conceive $\infty$ as magnitude, it nevertheless must
let $\infty$ figure as magnitude.}. Finite
% --- p. 47 ---
\opage{47}\setcounter{footnote}{0}%
these tangents surely cannot be, since the limit of finitude
has already been passed; infinite they cannot be either,
since the infinite has not yet been reached; but if they are neither
finite nor infinite, will not such tangents without
finitude and without infinity become just as interesting
as Baggesen's knives without handle and without blade? --- Had
Prof. St. only not had so downright rock-firm a faith in
``self-assumed or imagined authority,'' as he has revealed in these
% --- p. 48 ---
\opage{48}\setcounter{footnote}{0}%
corrections: then he would surely, instead of
arguing from an equation whose whole significance in this
context rests on the triviality that $\infty - a = \infty$,
to ``an infinitely great error,'' have thought a little, and reflected
on what the disputes known from the history of mathematics about the method of fluxions, the infinitesimal method and
the method of limits may well have to signify; and if
Prof. St. had come to think more closely about the history of
his science, he would not only have looked a little more
carefully at the Ramus --- and, alas, also Lagrange,
``quite peculiar proof,'' but surely also have thought
twice before so publicly and unreservedly
delivering a ``confession of faith'' like this: ``There is in the
first place nothing that can be called `the problem of infinity' in mathematics.''

How Prof. St.'s corrections justify the
parceling out by which he, in order to ``group the material,'' has
broken all connection and train of thought, even ``the less initiated'' may now well be able to judge. To be sure,
Prof. St. has, especially in the beginning, a notion
that a distinction must be made between essential and inessential errors, and in particular it seems to have become clear to
him that ``unmistakable misprints and slips of the pen'' must be counted among the inessential ones; but how little he in his great
zeal for precision is inclined to let any error be excused
by its being only a misprint, can be seen among other things
from the following example: ``On page 65 (Contr. 46) it is asserted that
$\frac{d^{2}y}{dx^{2}}$ is the increment with which $\mathrm{tg}\,\alpha$ varies when $x$ grows.
That is incorrect, for this increment is $\frac{d^{2}y}{dx}$, and here there is no misprint, as the whole development and the repetition of the expression
prove.'' Already in the expression ``is asserted'' there lies an arbitrary exaggeration, so that this expression is to that extent incorrect.
What is said about ``the whole development'' is likewise incorrect;
% --- p. 49 ---
\opage{49}\setcounter{footnote}{0}%
for the expression: $\frac{d^{2}y}{dx^{2}} = \frac{d \cdot \frac{dy}{dx}}{dx} = \frac{d\,\mathrm{tg}\,\alpha}{dx}$ shows surely
as clearly as possible what is to be understood by the
increment with which $\mathrm{tg}\,\alpha$, namely measured by $dx$, varies
when $x$ grows. What is said about ``the repetition of the expression''
is further incorrect, for $\frac{d^{2}y}{dx^{2}}$ occurs only a single
time in such a way that, by absolutely refusing to supply the
$dx$ with which the increment is measured, and by which it is thus first
completely determined, one can twist it into a misunderstanding.
That the repetition of an erroneous expression, even if such a
repetition had taken place, should without further ado be a
proof against the possibility of a misprint is, finally, a
mode of inference that Prof. St., in spite of his extraordinary
precision, nevertheless ought, if not for the sake of others, then at least for
his own, to be somewhat cautious about applying.

In Contr. (p. 34) it is, after all, expressly repeated that the remainder
which is to be added after $x^{n}$ is $\frac{x^{n}}{1 - x}$\footnote{``In the mathematical language the predicate is congruent with the subject''; from this proposition Prof. St. wants to prove that one does not
dare write $\frac{1}{1 - x} = 1 + x + x^{2} + \dots$ If it were hereby
said that by taking more and more terms in the infinite
series one constantly approached more closely the finite
% (the note runs on to p. 50)
value expressed by $\frac{1}{1 - x}$: then for $x > 1$ there would be ground here to raise an objection; but of such an approximation there is here no question. The sign of equality means only that the series: $1 + x + x^{2} + ....$, in order to correspond in the form of infinity to, to be congruent with, $\frac{1}{1 - x}$
in the form of finitude, must have an infinite number of terms; which
can also be expressed thus, that $\frac{1}{1 - x}$ by division produces
$1 + x + x^{2} + ....$, the division being continued without interruption. When there is now raised against this the schoolmasterish objection that
from $\frac{1}{1 - x}$ there is produced by division $1 + x + x^{2} + . .$
$x^{n} + \frac{x^{n+1}}{1 - x}$, even if n should be infinite: then this
objection is only a new proof of the thoughtlessness to which
the mechanics of the method of leaps can lead a mathematical leaper.
If ``development of concepts'' is in question here, it must surely
be permitted to ask about the difference between the conceivable and the inconceivable. But how is it % print: dct (misprint)
conceivable that after $x^{n}$ there can follow
an $x^{n+1}$, when n is to be infinite? What does it mean that
after $x^{n}$ a remainder $\frac{x^{n+1}}{1 - x}$ is added? It means indeed that the division,
instead of being continued, on the contrary is interrupted; but by the
division's being interrupted and the remainder added the infinite series has indeed
become finite. That a division is to be continued ad infinitum
thus becomes, by a mathematical decree: ``Remainder added!'' --- a decree
at the blackboard, determined so that it is to be interrupted somewhere in the
finite. If one is now prepared to endure the contradiction in so crude a
form, so long as the machinery of operations runs, and the calculation gives
the right result: then one can indeed readily scorn
``dialectic''; but one ought then at the same time to be
so consistent as to scorn every ``development of concepts.''}. What if now
someone, on the occasion of the fact that the same error $\frac{x^{n}}{1 - x}$ is undeniably
repeated, were by due trumpet blasts to proclaim that
Prof. St., one of ``the mathematicians of Europe,'' still
displays so gross an ignorance that he assumes
the remainder following $x^{n}$ to be $\frac{x^{n}}{1 - x}$!

Prof. St.'s ``Contribution'' bears, in the whole as in the
particular, the stamp of being a rush job occasioned by a certain
% --- p. 50 ---
\opage{50}\setcounter{footnote}{0}%
mood of the moment. To him who can now
grasp what Prof. St. calls ``essential errors'' and
``gross errors,'' and how he goes about correcting these
``essential errors,'' these ``gross errors,'' it will undoubtedly make
an imposing, not to say intimidating, impression to
hear Prof. St. conclude his 5th section with the following tirade
about Prof. Jürgensen: ``Can perhaps Prof. Jürgensen
% --- p. 51 ---
\opage{51}\setcounter{footnote}{0}%
show that the mathematical criticisms I find myself here compelled to make are unjustified or inessential? It
would undeniably be a surprise to me if he were to
attempt it, but succeed he cannot.'' --- To use
such language against such a man, one must be
conscious of one's superiority, and Prof. Steen is precisely a
man who is conscious of his superiority. Such a
``mathematician of Europe,'' who is conscious of his superiority, is of course no \textit{miles gloriosus}, but an imposing, intimidating figure. It is therefore to be expected that
this magnitude as negative as positive, this figure as
schoolmasterish as triumphant, will in the future,
as hitherto, maintain its place, show its superiority, and preserve its bearing, always imposing, intimidating, schoolmasterish, triumphant.

In all sciences, and then notably in the most exact
of all sciences, what matters is to be precise. But the
true demand for precision must always refer to a
difference between the essential and the inessential, and this difference in turn to a decisive point of view, a train of thought determining the whole as well as the particular.

Precision and completeness belong among the
categories whose dialectical double-sidedness shows itself especially in the determination of the relation of limits between philosophy and
the special sciences. What has been said in that respect about
the categories: obscurity, clarity and distinctness (Phil.
Propaed. pp. 108---109), then holds also of precision and completeness. ``Every scientific presentation of deeper content always has a multitude of luminous
points in the foreground, a multitude of obscure points in the background, and between both manifold transitions. Where
several sciences treat the same thing, each will make its own
distinct; but they will never make the same thing distinct in the
same way.'' The same with regard to precision.
What a development, according to its purpose, must determine with
% --- p. 52 ---
\opage{52}\setcounter{footnote}{0}%
the greatest precision, another may, according to its purpose, merely
suggest in fleeting, hovering strokes, just because it is to let
the emphasis of precision fall on other points.

If one abstracts from the considerations of clarity, distinctness, precision and completeness determined by the mutual
difference of the sciences, there remain only aphoristic notions of a distinct and an indistinct, a precise and
an imprecise, i.e. a correct and an incorrect.
What is then not correct is thus incorrect, and what is
not incorrect is correct. The aphoristic ``correct or
incorrect?'' becomes by abstraction a game, a game entertaining for pettifoggers, pedants and cavilers, a
literary one: ``odd or even?'' 

\begin{center}\rule{3em}{0.4pt}\end{center}

% --- p. 53 ---
\opage{53}\setcounter{footnote}{0}%

\begin{center}Misprints and Corrections.\end{center}

\begin{center}\rule{3em}{0.4pt}\end{center}

\begin{center}I. Philosophy and Mathematics.\end{center}

\noindent p. 10 line 11 from top $a - x$ read: $x - a$.

\noindent - - - 15---20 delete.

\noindent - 30 - 7 from bottom $\frac{A_{2}A_{1}}{1 + 1}$ read: $\frac{A_{1}A_{1}}{1 + 1}$.

\noindent - 36 - 5 - - $-\frac{x^{2}}{1.2}$ - $-\frac{x^{2}}{2}$.

\noindent - - - 4 - - $\frac{x^{3}}{1.2.3} - \frac{x^{4}}{1.2.3.4}$ read: $\frac{x^{3}}{3} - \frac{x^{4}}{4}$.

\noindent - - - 2 - - $\log \left( \frac{(1 + x)}{\log e} \right)^{\frac{2y^{2}}{1.2}}$ read:
\[ \left( \frac{\log (1 + x)}{\log e} \right)^{\frac{2y^{2}}{1.2}}. \]

\noindent - 37 - 8 from top $-\frac{x^{6}}{1.2.3.4.5.6}$ read:
\[ -\frac{x^{7}}{1.2.3.4.5.6.7}. \]

\noindent - 40 - 3 from bottom $l (-a)$ have fælleds med $l. a$

\noindent read: $2 l (-a)$ have fælleds med $l (a^{2})$.

\noindent - 43 - 11---28 from top delete.

\noindent - 44 - 2 and 11 from top add the factors $h^{m} k^{n}$.

\noindent - 49 after line 3 from top add: idet her endnu foreløbig ligesaalidt reflecteres paa nogen Forskjel

\noindent imellem $\infty$ og $\infty$, som paa nogen Forskjel

\noindent imellem $0$ og $0$.

% --- p. 54 ---
\opage{54}\setcounter{footnote}{0}%
p. 57 line 3 from top $f''''(x)$, . . read: $f'''(x)$ and $d^{3}y =$
\[ f'''(x)\ dx^{3}. \]

\noindent - 77 - 15 from top

\noindent $A = ds$, den uendelig read: $AB = ds$, det uendelig

\noindent lille Sector. lille Segment.

\noindent - - line 2 from bottom $\frac{dx}{dy}$ read: $\frac{dy}{dx}$.

\noindent - 82 - 16 - - et vist Forhold read: en vis Relation.

\noindent - - - 14 - - Forhold read: Relation.

\noindent - 83 - 5 and 10 from bottom $r^{2} \sin \Theta dr d\Theta d\alpha$ read:
\[ r^{2} dr \sin \alpha d\alpha d\Theta. \]

\noindent - 84 - 8 from top likewise.

\begin{center}II. Philosophical Propaedeutics.\end{center}

\noindent p. 92 line 16 from top Systemer --- det sphæroedriske \dots{}

\noindent read: Systemer, som, for at

\noindent nævne et Exempel, --- det sphæroedriske \dots{}

\noindent - 106 - 7 - - største, Teluret den mindste read:

\noindent mindste, Teluret den største.

\noindent - 139 - 12 - - her som Bindeled \dots{} read: her,

\noindent for at minde om et moderne

\noindent Synspunkt, som Bindeled.

\noindent - 161 - 6 - - den baconiske \dots{} read: den af

\noindent Apelt anførte baconiske \dots{}

\begin{center}III. Mathematics and Dialectic.\end{center}

\noindent p. 67 line 7 from bottom After $\frac{1.2.3 \dots n}{(a+1) \dots (a+n)}$

\noindent drops out $+ ....$

\noindent - 74 note line 2 from top $(1+x)^{1+x^{2}}$ read:
\[ (1+x)^{1-x^{2}}. \]

% --- p. 55 ---
\opage{55}\setcounter{footnote}{0}%
p. 77 line 2 from bottom Differentialcoifficienten read:

\noindent Differentialcoefficienten.

\noindent - 86 - 7 - - $\frac{du}{dx_{2}}$ read: $\frac{du}{dx_{2}} h_{2}$.

\noindent - 88 - 6 from top $\varphi'(x)\frac{n}{1} ....$ read: $\varphi'(x)\frac{h}{1}$.

\noindent - 95 - 3 - - af høiere end $r^{te}$ Orden \dots{}

\noindent read: af $r^{te}$ og høiere Orden \dots{}

\noindent - 96 - 9 - - $f''(x)\frac{dx}{1} + f'''(x)\frac{dx^{2}}{1.2}$ read:
\[ f''(x)\frac{dx}{1.2} + f'''(x)\frac{dx^{2}}{1.2.3}. \]

\noindent - - - 14 - - $f'''(x)\frac{dx^{2}}{1.2}$ \dots{} read:

\noindent $f'''(x)\frac{dx^{3}}{1.2.3}$ \dots{}

\noindent - 106 --- 1st line of text from bottom $\int_{x=b}^{x=b}(2xdx + dx^{2})$ \dots{} read:
\[ \int_{x=a}^{x=b}(2xdx + dx^{2}) \]

\noindent - 108 line 5 from top $\int_{0}^{x} 3x^{3}dx$ . . read: $\int_{0}^{x} 3x^{2}dx$ \dots{}

\noindent - - - 7 - - egentlige read: legemlige.

\noindent - 132 note line 3 from bottom παραλογιζόμενος read: παραλογιζόμενος.

\noindent - 138 - - 7 from top καθ, αὐτό read: καθ᾽ αὐτο. 

\noindent - - - - - - - αεί  read ἀει \dots{}

\noindent - 139 - - 8 - - \textit{Physie} read: \textit{Physic}.

\end{document}
